\documentclass[12pt]{article}

% Idioma e codificação
\usepackage[brazil]{babel}
\usepackage[utf8]{inputenc}
\usepackage[T1]{fontenc}

% Matemática e listas
\usepackage{amsmath, amssymb}
\usepackage{enumitem}

% Layout e cores
\usepackage{geometry}
\usepackage{graphicx}
\usepackage{xcolor}

% Cores institucionais
\definecolor{maincolor}{RGB}{0,100,50}
\definecolor{forestgreen}{RGB}{0,100,50}

% Cabeçalho e rodapé
\usepackage{fancyhdr}

% Blocos
\usepackage[most]{tcolorbox}
\newtcolorbox{block}[1]{%
  colback=green!5,
  colframe=maincolor,
  coltitle=white,
  title=#1,
  fonttitle=\bfseries,
  boxrule=1.5pt,
  arc=3mm,
  left=6pt,
  right=6pt,
  top=6pt,
  bottom=6pt
}

% Margens
\geometry{margin=2.5cm}

% Fancyhdr
\pagestyle{fancy}
\fancyhf{}

\renewcommand{\headrulewidth}{2pt}
\renewcommand{\footrulewidth}{2pt}

\renewcommand{\headrule}{\hbox to\headwidth{%
  \color{forestgreen}\leaders\hrule height \headrulewidth\hfill}}

\renewcommand{\footrule}{\hbox to\headwidth{%
  \color{forestgreen}\leaders\hrule height \footrulewidth\hfill}}

% Altura do cabeçalho (corrigida)
\setlength{\headheight}{52pt}
\addtolength{\topmargin}{-20pt}

% Cabeçalho
\fancyhead[L]{%
  \includegraphics[height=1.2cm]{Universidade_de_Rio_Verde.jpg}
}

\fancyhead[R]{%
  \raggedleft
  \textcolor{forestgreen}{%
    \small
    \textbf{Trabalho Prático}\\
    Me.~Sergio Souza Novak\\
    \textit{Processo de Software}
    }
    }
    
    % Rodapé
    \fancyfoot[C]{\textcolor{forestgreen}{\thepage}}
    
    % =========================
    \begin{document}
    % =========================
    
    \vspace*{0.2cm}
    
\begin{center}
  {\Large \textbf{Trabalho Prático: Documentação para Registro de Software no INPI}}\\[0.2cm]
\end{center}

\vspace{0.4cm}

\textbf{Objetivo:} Realizar a preparação e organização da documentação necessária para o registro de um software no INPI (Instituto Nacional da Propriedade Industrial).

\textbf{Instruções:}
\begin{itemize}
    \item O trabalho deve ser realizado em \textbf{duplas}.
    \item \textcolor{red}{\textbf{Os dois da dupla devem entregar o trabalho.}}
    \item Os membros da dupla devem constar como os \textbf{autores} do software na documentação.
    \item Os alunos podem utilizar qualquer software desenvolvido por eles ao longo da graduação, independentemente da sua complexidade.
    \item É obrigatório o preenchimento de todos os documentos (.odt) fornecidos.
    \item Deve ser gravado e entregue um \textbf{vídeo explicativo} demonstrando o funcionamento do software.
    \item A entrega final deverá ser feita através de um único arquivo compactado (ZIP).
    \item \textcolor{red}{\textbf{Ambos os alunos da dupla devem realizar a entrega do trabalho no sistema.}}
\end{itemize}

\textbf{Estrutura do Arquivo ZIP:}

O arquivo compactado deve obrigatoriamente conter a seguinte estrutura de diretórios e arquivos:

\begin{verbatim}
.
|-- Autores
|   |-- [Nome do Aluno 1].txt
|   \-- [Nome do Aluno 2].txt
|-- Código
|   \-- app.py
|-- comprovação de chave SHA256
|   \-- Imagem colada.png
|-- DECLARAÇÃO DE AUTENTICIDADE.odt
|-- FORMULARIO.odt
|-- imagens
|   \-- Imagem colada.png
|-- TERMO DE PARTICIPAÇÃO DOS AUTORES.odt
\-- vídeo.mp4
\end{verbatim}

\vspace{0.4cm}

\textbf{Observações Finais:}
Certifiquem-se de que todos os formulários e termos (.odt ou docx) estejam devidamente preenchidos com os dados da dupla (autores). A comprovação da chave SHA256 deve corresponder ao código-fonte (app.py ou outro arquivo único, podendo ser zipado) e documentação fornecidos.

% =========================
\end{document}
% =========================
